% !TeX root = ../papers/nc-p3.tex
\section{二つのメビウス成分と一次元点グラフの全被覆}
\label{sec:split-binary-pencil-charts}
\begin{theorem}\label{thm:split-binary-pencil-charts}
前節の三次残余で $E$ が二つの相異なる $(1,1)$ グラフの和なら、全ての窓は標準Koszul左右AS正則である。
\end{theorem}
\begin{proof}[二つの節点での座標]
反対二位置のある二点で $f$ が零行列なら四テンソルもそこで消え、系\ref{cor:binary-vertical-ternary}の
三支持の議論で済む。以下、隣接・反対のどの二位置にも零行列がないとする。
二成分の交点は長さ2である。まず二単点の場合を扱う。
どの位置でもその二点の像は相異なる。四変数の各々でこの二点に対応する座標を $0,\infty$ とする。
ここで四変数のラベルは一つの特異点での核を用いて合わせる。実際、最初の二変数の
行列式の特異点では、その両核を残りの二変数に入れた四組が $f$ の全ての一階偏微分を消す。
これは階数1行列の行列式の微分が左右核での評価に比例するためである。
この四組から任意の隣接二座標を取っても同じ特異点の移動となる。
従って二つの四組は全て0と全て無限遠になり、$f$ は二変数ずつの積六項だけを持つ。
係数が一つでも零ならその相補二位置の $(0,0)$ で零行列となる。
各座標の非零定数倍で相補係数を等しくできる（三対の比の文字は独立で、代数閉体上そのトーラス写像は全射）。
ベクトル座標を $u,v,z,t$ とすると
\[
 F=a(uv+zt)+b(uz+vt)+c(ut+vz),\qquad abc\ne0 .
\]
ここでの四つの座標は独立に選んだもので、点グラフ $C$ に沿った同一視とは一般に異なる。
$u,v$ を同じ $C$ に戻すと二点は $P_0 u,P_1 v$ となる。
$z,t$ も同様に $P_2 z,P_3 t$ と書く。$\Theta=P_0^{-1}P_1,\ \eta=P_2^{-1}P_3$ と置く。

$u,v$ の二成分は $v=\lambda u,\lambda^{-1}u$ で、
$bc(\lambda+\lambda^{-1})=a^2-b^2-c^2$, $\lambda^2\ne1$ である。
$m=-(c+b\lambda)/a,\ n=-(b+c\lambda)/a$ とすると $mn=\lambda$ で、
第一成分での残る二根は $z=mu,t=nu$、第二成分では $m,n$ を逆数にする。
特に $m^2\ne1$ である。同じ $u,z=mu$ で二つの $v$ が核になるなら反対二位置で零行列になるからである。
補題\ref{lem:cubic-secant-isomorphism}の辺で、第一成分は
\[
 \{P_0 x,P_1(\lambda x)\}\ \longmapsto \{P_2(x/m),P_3(\lambda x/m)\},
\]
第二成分は $\lambda,m$ を逆数にした対応を与える。実際、$t=nx,u=x$ を固定したとき
$v$ 全体で消える根は $z=x/m$ である。これは $F$ に代入して従う。
四移動の継ぎ目があれば補題と同じく終点の二座標を一緒に移すので相対変換 $\eta$ には影響しない。
\end{proof}

\begin{proof}[混合した変換の制限]
$\Lambda(x)=\lambda x$ とし、二つの変換 $\Theta\Lambda,\Theta\Lambda^{-1}$ を調べる。
どちらかが恒等なら $\Theta$ は対角である。ともに非自明な対合なら二つの行列の跡が消えて
$\Theta$ は反対角である。どちらも対合でも恒等でもなければ、順序を忘れた両像は二次曲線で
同じ三次式に入るため一致する。順序を戻す二つのグラフは互いの逆であり
$(\Theta\Lambda^{-1})^2=\Lambda^{-2}$ が射影的に成り立つ。
非自明な倍率の平方根はその二固定点を保ち、反対角行列の平方は恒等だから $\Theta$ は対角である。

片方だけが非自明な対合で他方が恒等でない場合、第一成分を対合としてよい。
$\Theta=\bigl(\begin{smallmatrix}A&B\\C&-\lambda A\end{smallmatrix}\bigr)$ と書くと $A\ne0$。
第一成分上の割線因子は次数2で直線を被覆し、被覆対合は $\Theta\Lambda$ である。
辺が同型なので、$M(x)=mx$ とすると終点の対合は
$\eta\Lambda=M^{-1}\Theta\Lambda M$、ゆえに $\eta=M^{-1}\Theta M$ となる。

命題\ref{prop:twisted-cubic-secant-cubic}の二因子のうち、残る二根の一致を検出する双一次式を $H$ と書く。
その零点は、第二成分の逆グラフを含む。実際、順序を忘れた三次式は第一成分の直線と第二成分の二次曲線を含み、
後者の引戻しは元のグラフと逆グラフの和である。逆グラフは初めの二成分と異なるので $H$ に入る。
双一次式の次数により両者の式は比例する。残る二根の一致は $z=\eta t$ であるから、
$F$ の四係数にこの交代縮約を適用し、共通の非零倍を除くと
\[
 H=m Cuv-nAu-mAv-B/m .
\]
逆グラフは $u=(\Theta\Lambda^{-1}\Theta)v$ であり、その式を $G$ とすると
\[
 G=AC(\lambda^{-1}-\lambda)uv+
 (BC/\lambda+A^2\lambda^2)u-(A^2/\lambda+BC)v-AB(\lambda^{-1}-\lambda).
\]
$B,C$ が共に非零なら $uv$ と定数の係数比で $m^2=1$ となり不可能である。
$B\ne0,C=0$ なら $u,v$ の係数比で $m^2=-1/\lambda^2$、
$v$ と定数の比で $m^2(1-\lambda^2)=1$ となり矛盾する。
$C\ne0,B=0$ は四座標を一緒に逆数にして同じ場合に帰着する。
混合型でも $\Theta$ は対角である。

従って全ての隣接位置で二つの節点の像の非順序対は同じである。これを $0,\infty$ と正規化する。
全ての位置の割線三次曲線は、$J:x\mapsto-x$ の点対 $\{x,-x\}$ の直線 $L$ を含む。
実際、$\Theta$ が対角なら元の二成分は倍率変換のグラフである。
一方が $-1$ の倍率ならすでに $L$ がある。そうでなければ二次像は一致し、二つの非自明な倍率は互いの逆となる。
各成分から割線因子への次数は1で、終点でも同じだから $\eta$ は対角である
（反対角なら両成分が直線への二重被覆になる）。
$\Theta$ が反対角なら元の二グラフはともに $x\mapsto k/x$ 型の対合で、二つの相異なる直線を与える。
各成分での次数2により終点でも二つの相異なる対合が必要で $\eta$ は反対角である。

いずれの場合も残る縮約 $H$ は共通の物理座標 $p,q$ において、
斉次化すれば $h_1p+h_2q$ の形である。対角の場合は $F$ の $z,t$ の一次係数を使う。
反対角の場合は $uv$ と定数を使い、第二物理座標が $v$ の逆数の定数倍である。
原点と無限遠の同時交換でもこの形は保たれる。$H$ が恒等零なら割線三次式も恒等零となり不可能である。
一方の係数だけ非零なら、その二本の縦横の零線は順序交換で保たれず、その交換像は
元の二グラフにも入らないので三次式の対称性に反する。
対角で逆グラフどうしの場合、二次曲線の残余直線の像は $H$ から来なければならない。
反対角の場合も $H$ の倍率グラフは元の二つの対合と異なる。
よってこの倍率グラフは直線に写らなければならず、倍率は $-1$ である
（恒等なら対角二次曲線、その他の非対合も二次曲線に写る）。これで $L$ の包含を得た。

四位置の節点座標での同時 $J$ は六項の $F$ を保ち、物理座標でも同じ $J$ なので角括弧も保つ（射影的に）。
四移動も二節点の像の対を保つので $J$ と可換である。
従って割線三次の同型は $J$ と可換で、固定軌跡の正次元成分 $L$ は次の $L$ に写る。
実際、対称平方平面の固定点集合は $L$ と点 $\{0,\infty\}$ である。
これで両方向の割線列ができ、全次数正則性を得る。
\end{proof}

\begin{proof}[一つの二重交点]
二成分が一点で接するとする。各成分を第一射影でパラメータ化し、接点を無限遠に取る。
各位置で第一成分の前向きの座標を $x$ 自身に取ると、第二成分は接点の一次ジェットを共有するので、
四つの座標は $x,x+e,x+f,x+g$ となる。$e\ne0,f-e\ne0$ は隣り合うグラフの相異性による。
$z=v$ と $z=v+f-e$ 上で、$u,t$ の行列の二根は前後の座標である。
従って四テンソルの商 $F$ の制限は非零定数倍の $(u-v)(t-v)$ と $(u-v+e)(t-v+e-g)$ である。
倍率が定数なのは斉次化した同じ次数の切断であり、零行列もないからである。
$z$ で一次補間し、$v^3,v^2$ の係数を消すと二定数は等しく $g=f+e$ で、
\[
 F=(u-v)(t-v)+\frac{z-v}{f-e}\bigl(-f(u-v)+e(t-v)-ef\bigr)
\]
となる。第一成分の二根は $v=u$ で $z=u,t=u+f$ であり、
$t=u+f,u$ の組から来る逆側の根も $z'=u+f$ である。
従って第一成分の割線の移動は、始点の組 $(x,x)$ から終点の組 $(x+f,x+f)$ を使う。

物理座標の相対変換を同じく $\Theta,\eta$ とすると、始点の二グラフは
$\Theta$ と $\Theta T_e$（$T_e(x)=x+e$）である。
いずれかが恒等なら $\Theta$ は無限遠を固定する。ともに対合なら二つの跡の差が
$e\Theta_{21}$ なのでやはり固定する。ともに対合でも恒等でもなければ二次曲線の一致により
互いに逆で、$\Theta^2=T_{-e}$ だから唯一の固定点の無限遠を保つ。
混合なら対合の成分を第一に選ぶ。無限遠が固定されないと仮定し、四座標を一緒に平行移動して
$\Theta(x)=k/x,\ k\ne0$ とできる。第一成分の対合の移動から
$\eta=T_f\Theta T_{-f}$ である。残余縮約 $H$ は先と同じく第二成分の逆グラフを含まねばならない。
このグラフの式は $u=(\Theta T_e\Theta)v$、すなわち $euv+ku-kv=0$ である。
その上の $(u,v)=(0,0)$ では $F$ の二根は $(z,t)=(0,f)$ だが $\eta(f)=\infty$ で一致しない。
$H$ の消滅に反する。従って全辺で無限遠の像が一致し、接点を運ぶと物理座標が一定の $p$ になる。
全位置で $D_j=2p$ を使えば前節同様に両方向の割線列を得る。
\end{proof}

\begin{corollary}\label{cor:all-one-dimensional-point-graphs}
両側点ファイバーで一次元の点グラフを持つ枝では、全標準窓が有限の正則性・添字保持同型判定で被覆される。
純部分が空間有理正規三次となる最後の枝では、有限Frobenius検査と点グラフ検査を通った候補は、
埋込み次元2以下なら全て正則である。埋込み次元3では定理\ref{thm:maximal-tangent-ore}の判定を用いる。
三次曲線が唯一の一次元の台であることを課し、四つの完全三次系で座標をそろえると、この枝の
パラメータは四移動 $S$ と多重一次形式の射影類 $[f]$ である。許容域は四括弧の積と $f$ から作る
四テンソルについて上記の有限検査を課したものである。
同型条件は同時 $\mathrm{PGL}_2$ 作用による $S$ の共役と $[f]$ の対応であり、
周期は一・二位置回転について同じ同型条件を試し、いずれも失敗すれば4となる。
\end{corollary}
\begin{proof}
系\ref{cor:point-curve-high-coverage}、定理\ref{thm:elliptic-binary-cubic-pencils}、
系\ref{cor:binary-vertical-ternary}、定理\ref{thm:doubled-binary-pencils}と本定理を合わせる。
同型は唯一の一次元の被約曲線を保つ。最初の三辺の点列がそこで対角であるため
四位置の基底変更はこの曲線上で同じ射影変換を誘導する。完全三次系なので各周囲空間の射影変換もこれで決まる。
第四辺で $S$ が対応し、四テンソルの交叉直線から $[f]$ が対応する。逆も各関係の復元により従う。
標準四周期構造を忘れた同型と周期については定理\ref{thm:canonical-four}と系\ref{cor:period-orbits}を用いる。
\end{proof}
